\chapter{The Command Line Interface}
\label{ch:cli}

\newthought{The command line is where the design becomes visible.} \Jtwo\ deliberately
exposes one executable, \cli{Jarvis}, and asks you to name \emph{one intent} per
invocation --- and the intents are not an arbitrary menu. Three of them are graded by how
much of your task card they are willing to execute, from reading it to running it; three
more act on scans already running, and never touch a card at all. Learn that shape and the
\textsc{cli} stops being a list of flags to memorise: you will know which command to reach
for, what it started on your machine, what it deliberately did not start, and how to read
what it printed back.

\section{One binary, one intent, one card}
\label{sec:cli-shape}

Every invocation has the same shape:

\begin{terminal}
\begin{jcmd}
Jarvis <intent> <card>.yaml [flags]
\end{jcmd}
\end{terminal}

The intent comes first because it is the decision you actually make; the card is the
noun it acts on. \cli{Jarvis -h} prints the whole surface, and it does not print a flat
list --- the commands arrive already sorted into the five groups of
Table~\ref{tab:cli-groups}, which are the groups you actually work in:

\begin{table}[ht]
  \footnotesize
  \begingroup
  \setlength{\tabcolsep}{0pt}
  \begin{tabular}{@{}p{0.255\linewidth}p{0.213\linewidth}p{0.532\linewidth}@{}}
    \toprule
    \textbf{Group} & \textbf{Commands} & \textbf{What the group is for} \\
    \midrule
    Scan workflow &
      \clih{validate} \clih{check} \clih{run} &
      the same card, executed at three depths (\S\ref{sec:cli-ladder}) \\
    Data export &
      \clih{convert} &
      refresh the \textsc{csv} snapshots of a run's \textsc{hdf5} \\
    Plots &
      \clih{gen-plot-yaml} \clih{plot} &
      generate a plot card from a finished run, then render it \\
    Runtime control &
      \clih{monitor} \clih{ps} \clih{kill} &
      look at, or end, scans running right now --- no card involved \\
    Extensions \& projects &
      \clih{project} \clih{portal} \clih{operas} &
      doors into the neighbouring tools \\
    \bottomrule
  \end{tabular}
  \endgroup
  \caption[The five command groups printed by \clih{Jarvis -h}.]{The five command
  groups printed by \clih{Jarvis -h}. Within the first group the order is the order you
  should use them in.}
  \label{tab:cli-groups}
\end{table}

Those group names are the CLI's own, not this book's, and they are worth trusting: when a
command moves or appears, the help panel is where it shows up first. Two habits follow.
Run \cli{Jarvis -h} when you have forgotten \emph{which} command, and
\cli{Jarvis <command> -h} when you have forgotten \emph{which flag} --- the help text
says so itself, and every command accepts both \cli{-h} and \cli{--help}.

\paragraph{The card never changes between intents.} This is the single most important
consequence of the design, and it is worth stating as a rule: the file you hand to
\cli{Jarvis validate} is byte-for-byte the file you hand to \cli{Jarvis run}. There is
no ``test card'' and no debug variant to keep in sync. Gating a card in continuous
integration therefore gates \emph{the card you will run}, and a card that survives
\cli{check} needs no edit before production --- only more points.

\paragraph{The CLI refuses to guess.} Called with nothing, it says so and stops rather
than defaulting to some plausible action:

\begin{terminal}
\begin{jcmd}
Jarvis
\end{jcmd}
\termsep
\begin{jout}
Provide a task YAML or a subcommand
(run|check|validate|convert|monitor|plot|portal|operas|project). See Jarvis -h.
\end{jout}
\end{terminal}

The same instinct governs conflicts. Naming two intents at once --- a subcommand plus a
legacy flag that means something different (Section~\ref{sec:cli-legacy}) --- is a usage
error, not a precedence puzzle to be resolved silently in your favour.

\section{The ladder of commitment}
\label{sec:cli-ladder}

The three commands in the \emph{Scan workflow} group run the same card at three depths.
Almost every practical question about ``which command do I use'' is a question about
which rung you are on --- which is why the help panel lists them in this order rather
than alphabetically. Table~\ref{tab:cli-ladder} is that ladder.

\begin{table}[ht]
  \footnotesize
  \begingroup
  \setlength{\tabcolsep}{0pt}
  \begin{tabular}{@{}p{0.255\linewidth}p{0.213\linewidth}p{0.532\linewidth}@{}}
    \toprule
    \textbf{Intent} & \textbf{Cost} & \textbf{What it starts, and what it proves} \\
    \midrule
    \clih{validate} & seconds &
      Nothing. Reads and checks the card against the \textsc{yaml} contract. Proves the
      card is \emph{sayable}. \\
    \clih{check}    & minutes &
      Redis, Workers, your real calculators --- on about ten points. Proves the card is
      \emph{runnable}: paths resolve, programs execute, outputs parse, the likelihood
      evaluates. \\
    \clih{run}      & hours &
      Everything, for as long as the sampler asks. Proves nothing new about the card ---
      it produces physics. \\
    \bottomrule
  \end{tabular}
  \endgroup
  \caption{The three depths at which \clih{Jarvis} will execute one card.}
  \label{tab:cli-ladder}
\end{table}

The rungs are cheap in exactly the order you need them to be, and each one catches a
different class of mistake. \cli{validate} catches what you wrote: a misspelled
\yamlkey{Sampling.Method}, a key that no longer exists, a bound that contradicts a
distribution. \cli{check} catches what your \emph{machine} does with what you wrote: an
\yamlkey{entry} path that is correct on the cluster and wrong on the laptop, a program
that needs an environment variable you forgot to set, an output file whose format
changed between versions of an external code. No amount of validation finds those,
because none of them are in the card.

That is why the failure mode this ladder exists to prevent is so specific: discovering a
broken path forty minutes into a scan, after a sampler has already spent your evening.

\section{\clih{validate}: the cheapest gate}
\label{sec:cli-validate}

\begin{terminal}
\begin{jcmd}
Jarvis validate scan.yaml
Jarvis validate scan.yaml --strict   # also fail on warnings
Jarvis validate scan.yaml --json     # machine-readable report
\end{jcmd}
\end{terminal}

\cli{validate} loads a card and checks it against the full \textsc{yaml} contract ---
every block in Part~\ref{part:task-card}, plus the keys the chosen sampler declares for
itself --- without starting Redis, a Worker, or an Archiver. \cli{Jarvis run} and
\cli{Jarvis check} run this same gate before they do anything else, so \cli{validate}
is not an extra safety net bolted on the side. It is the first thing a run does, made
available on its own, for the moment \emph{before} you are willing to pay for a run.

A card that passes prints one line and exits \code{0}:

\begin{terminal}
\begin{jcmd}
Jarvis validate bin/Example_Dynesty_Operas_Quick.yaml
\end{jcmd}
\termsep
\begin{jout}
Config validation successful. The task YAML meets the V2 contract.
\end{jout}
\end{terminal}

A card that fails prints every problem it found, then exits \code{2}:

\begin{terminal}
\begin{jout}
Config validation failed (1 error(s), 0 warning(s)):
  [error] JV2-MTH-003  Sampling.Method
          'Dynastey' is not implemented in Jarvis-HEP V2. Available: AM, AMMCMC,
          AdaptiveBridson, Bridson, CSV, DEMCMC, DRAM, Dynesty, Ensemble,
          EnsembleMCMC, Grid, MCMC, MultiNest, PT, PTEnsemble, PTMCMC, Random
\end{jout}
\end{terminal}

Every finding carries the same four fields, and each one is there to save you a specific
step:

\begin{keylist}
  \item[{[error] / [warning]}] Severity. Errors fail the card; warnings do not, unless
        you asked for \cli{--strict}.
  \item[JV2-MTH-003] A stable code. It does not change between releases, so it is safe
        to grep for in \textsc{ci} logs and it is the lookup key in
        Appendix~\ref{app:troubleshooting}.
  \item[Sampling.Method] The exact dotted path in \emph{your} card. No line numbers, no
        guessing which of four \yamlkey{Method} keys it meant.
  \item[the message] What is wrong and --- wherever the answer is a closed set, as here
        --- what the legal values are.
\end{keylist}

\paragraph{\clih{--strict}.} Warnings mark things that will not stop a run but will
probably not do what you meant: a key that \Jtwo\ ignores, a dead
\yamlkey{Operas.call\_mode}, a setting the chosen sampler silently overrides.
\cli{--strict} promotes them to hard failures. Make it a habit before a long production
run: the card you are about to spend a night on is exactly the card whose ``it runs,
just not quite the way you think'' should be caught first.

\paragraph{\clih{--json}.} The same report, on standard output, as one \textsc{json}
object --- \code{ok}, the resolved \code{task\_yaml} path, the \code{scan\_name}, and an
\code{issues} array of the four fields above:

\begin{terminal}
\begin{jout}
{
  "issues": [
    {
      "code": "JV2-MTH-003",
      "hint": null,
      "level": "error",
      "message": "'Dynastey' is not implemented in Jarvis-HEP V2. ...",
      "path": "Sampling.Method"
    }
  ],
  "ok": false,
  "scan_name": "EggBox_Dynesty_Operas_Quick",
  "task_yaml": "/home/you/Eggbox/bin/Example_Dynesty_Operas_Quick.yaml"
}
\end{jout}
\end{terminal}

Note which stream each form uses: the human report goes to standard error, the
\textsc{json} report to standard output. A \textsc{ci} job can therefore pipe
\cli{--json} straight into a parser without filtering log noise out of it first. The
\sampler{MultiNest} and \sampler{Dynesty} chapters walk a real card through this gate
key by key (Section~\ref{sec:nested-validation}).

\section{\clih{check}: the workflow rehearsal}
\label{sec:cli-check}

\begin{terminal}
\begin{jcmd}
Jarvis check scan.yaml
Jarvis check scan.yaml --strict --console-level DEBUG
\end{jcmd}
\end{terminal}

\cli{check} is a full run in miniature. It starts the same Redis, the same Workers, and
the same calculators as production, then feeds them a handful of points --- fixed ones
from a \textsc{csv} if your card resolves to one, or about ten drawn from your real
\yamlkey{Sampling.Method} if it does not. Everything downstream is genuine: the same
shell commands, the same \Portal\ readers and writers, the same \Operas\ calls, the same
likelihood expression.

Two flags earn their keep here. \cli{--strict} applies to the same validation gate as
above. \cli{--console-level DEBUG} raises what reaches the screen, which is what you want
the first time a card fails on a new machine --- the files were already recording at
\code{DEBUG} either way (\S\ref{sec:cli-logging-flags}).
Chapter~\ref{ch:check-modules} covers where the points come from and how to read the
per-Sample output that follows.

\section{\clih{run}: the real thing}
\label{sec:cli-run}

\begin{terminal}
\begin{jcmd}
Jarvis run scan.yaml
Jarvis run scan.yaml --resume    # take the checkpoint, no prompt
\end{jcmd}
\end{terminal}

One command, in order: validate the card (clear messages and exit code~2 on a missing or
invalid file), start Redis bookkeeping, spawn the Workers and the Archiver, drive the
sampler until it says it is finished, then tear the whole family down --- including any
\code{Jarvis-Redis:<scan>} service \Jtwo\ started itself. If a checkpoint from an
interrupted run is sitting there, you are asked whether to resume before any of that
happens; \cli{--resume} answers yes non-interactively
(Chapter~\ref{ch:checkpoint}).

\subsection{The flags}

\begin{keylist}
  \item[--resume] Accept an existing checkpoint without prompting. The right flag for
        batch queues and \cli{nohup}, where there is no terminal to answer the question.
  \item[--strict] Promote validation warnings to errors, exactly as in
        Section~\ref{sec:cli-validate}. Costs nothing; run it this way.
  \item[--skip-draw-flowchart] Skip the \file{flowchart.json} and \file{flowchart.png}
        export. Useful on headless machines and in repeated short runs where you already
        know the workflow graph.
  \item[--console-level LEVEL] Minimum severity that reaches the screen: \code{DEBUG},
        \code{INFO}, \code{WARNING}, \code{ERROR}, or \code{CRITICAL} (default
        \code{WARNING}).
  \item[--debug, -d] Shorthand for \cli{--console-level DEBUG}.
  \item[--silence, -s] Turn console output off entirely. Files are still written, so this
        loses nothing --- it is the flag for a scan whose terminal you are about to stop
        watching.
\end{keylist}

\subsection{Deciding how much to watch}
\label{sec:cli-logging-flags}

\textbf{The files always record \code{DEBUG} and above.} That is not adjustable, and it is
the useful half of the design: you cannot start an overnight scan that fails at 02:00
having logged nothing worth reading.

So the only logging decision is how much of it reaches your screen, and the default answer
is \code{WARNING} --- a quiet terminal that speaks up when something is wrong. Raise it
with \cli{--console-level INFO} to watch progress, or \cli{--debug} when you are chasing a
specific failure. The file under \file{logs/<scan>/} is identical either way, so the
post-mortem recipe in Chapter~\ref{ch:logging} works no matter how the run was invoked.

\subsection{One scan per name}

Starting a run whose \yamlkey{Scan.name} already belongs to a live controller is refused
before Redis comes up:

\begin{terminal}
\begin{jout}
a Jarvis scan named 'EggBox_Dynesty' is already running; inspect it with
`Jarvis monitor R1` or `Jarvis ps R1` before starting another run
\end{jout}
\end{terminal}

Two scans sharing a name would share an output directory, a checkpoint file, and a set of
Redis keys, and the second one would quietly corrupt the first. The message hands you the
reference of the scan already holding the name so you can go and look at it
(\S\ref{sec:cli-ps-kill}) rather than guess.

\subsection{What it looks like on your machine}

Every managed process sets its own \textsc{os}-visible title, so \cli{ps} and \cli{top}
show you a scan instead of a column of identical \code{python3} entries:
\code{Jarvis:<scan>} for the control process,
\code{Jarvis-\allowbreak{}Worker-\allowbreak{}N:<scan>},
\code{Jarvis-\allowbreak{}Archiver:<scan>}, \code{Jarvis-\allowbreak{}FileOperation},
and \code{Jarvis-\allowbreak{}Redis:<scan>}. Section~\ref{sec:cli-ps-kill} puts that to
work.

\subsection{Stopping}

\textbf{Ctrl+C} (\textsc{sigint}) is the supported stop: Workers finish or abandon
cleanly, the checkpoint is saved, and children are reaped. \textbf{Ctrl+Z} is
deliberately ignored during a scan --- a suspended control process would leave Workers
and Redis half-alive, holding Redis keys and file handles that nothing is left to
release.

\subsection{The closing line}

A finished run prints one summary line to standard error before it exits:

\begin{terminal}
\begin{jout}
RunOutcome status=success submitted=20 completed=20 failed=0 archived=20 exit=0
\end{jout}
\end{terminal}

The counters are the same ones \file{run\_summary.json} records
(Chapter~\ref{ch:monitoring}), and \code{status} is the field the exit code is derived
from --- not the reverse. Section~\ref{sec:exit-codes} gives the mapping.

\section{The second terminal}
\label{sec:cli-second-terminal}

Here is a habit worth forming, and a design decision worth noticing behind it: while a
scan runs, everything you want to \emph{do} to it happens from a second terminal, and
none of it disturbs the first. \cli{Jarvis monitor} reads one status snapshot.
\cli{Jarvis ps} lists the processes. \cli{Jarvis convert} exports whatever data has
been flushed so far. All three are read-only with respect to the scan's state; you can
run them whenever you like, as often as you like, and stop.

That is possible because a scan's live state is not held inside the control process ---
it is in Redis and on disk (Chapter~\ref{ch:distributed-runtime}). Any tool that can
reach those can observe a scan without asking it to cooperate, which is also why none of
these commands need a \textsc{pid} or a socket.

\subsection{\clih{monitor}}

\begin{terminal}
\begin{jcmd}
Jarvis monitor       # which scans are running?
Jarvis monitor R1    # snapshot of that one
\end{jcmd}
\end{terminal}

Bare, \cli{monitor} lists the running scans and stops. Given a reference it attaches to
that scan's Redis and prints one read-only snapshot --- queue depths, Worker states,
throughput, resource use --- then exits. Chapter~\ref{ch:monitoring} reads a snapshot
field by field; Section~\ref{sec:cli-ps-kill} explains the reference.

\subsection{\clih{convert}: refreshing the \textsc{csv} snapshots}
\label{sec:cli-convert}

\begin{terminal}
\begin{jcmd}
Jarvis convert scan.yaml
\end{jcmd}
\end{terminal}

The Archiver already writes \file{DATABASE/samples.csv} at the end of every run
(Section~\ref{sec:database-csv}), so \cli{convert} has two jobs, and both are worth
knowing. It is the recovery path when that export never happened --- an interrupted run,
a deleted \textsc{csv}, or any other \file{*.hdf5} under \file{DATABASE/} that never got
one. And it is how you look at a long scan's results \emph{before} it finishes: convert
in the second terminal, open the \textsc{csv}, and decide whether the run is worth
letting continue.

It resolves \yamlkey{task\_result\_dir} from the card (path resolution only --- no Redis,
no full validation gate) and refreshes \emph{every} \file{*.hdf5} directly under
\file{DATABASE/} into a sibling \file{.csv} of the same name, one column per observable,
with nested values flattened to \textsc{json} text by the same rule \file{samples.csv}
itself uses.

\textbf{Run it as often as you like.} \cli{convert} compares each \textsc{hdf5} against
the \textsc{csv} beside it and rewrites only what actually changed, so repeated runs on a
growing scan are cheap and a run on a finished scan is a no-op. It is safe mid-scan for
the same reason it needs no locking: it only reads \file{DATABASE/*.hdf5} and writes
\textsc{csv} beside it, never touching the running scan's Redis state, so you get whatever
the Archiver has flushed so far.

Each file reports its own outcome: \code{converted} with a row count,
\code{skipped\_unchanged} when the \textsc{csv} already matches its source, \code{empty}
when the \textsc{hdf5} holds no records yet, or \code{missing} when it is unreadable or
absent. Exit code \code{0} if at least one file was converted or already up to date,
\code{1} if every file was empty or missing, or if no \file{*.hdf5} exists at all.

\subsection{\clih{gen-plot-yaml}: a plot card from a finished run}
\label{sec:cli-gen-plot-yaml}

\begin{terminal}
\begin{jcmd}
Jarvis gen-plot-yaml scan.yaml    # writes images/<scan>/plot.yaml, prints its path
Jarvis plot images/<scan>/plot.yaml

# or chain them: the generator prints the path it wrote
Jarvis plot "$(Jarvis gen-plot-yaml scan.yaml)"
\end{jcmd}
\end{terminal}

The two \emph{Plots} commands are deliberately separate steps rather than one
``\,plot my scan'' button. \cli{gen-plot-yaml} reads a finished run's outputs and writes
\file{images/<scan>/plot.yaml}, a \JPlot\ scene card seeded with that scan's own variables
and likelihood --- a starting point, not a final figure. You then edit it and render it
with \cli{plot}, as many times as the figure needs. Generation happens once; rendering
happens as often as you like, which is the asymmetry the split exists to serve
(Chapter~\ref{ch:jarvisplot}). On success the command prints the path it wrote --- the
last line above uses that to chain the two halves in one go.

\begin{jwarn}
\cli{gen-plot-yaml} rewrites \file{plot.yaml} in place, without asking. Once you start
tuning a figure, save your version under another name --- \file{plot\_paper.yaml} --- or
the next regeneration silently discards it.
\end{jwarn}

\section{Runtime control: pick a scan, then act on it}
\label{sec:cli-ps-kill}

A distributed scan is a small family of processes, and occasionally --- a killed
terminal, a lost \textsc{ssh} session, a control process that took a signal the wrong
way --- part of that family outlives the run. Worse, on a shared machine there may be
several families, and only one of them is yours to end. Guessing at \textsc{pid}s from
\cli{ps aux | grep python} is exactly the operation you should not have to do by hand.

All three \emph{Runtime control} commands therefore work the same two-step way: run them
bare to see what is running, then run them again naming one scan.

\subsection{Step one: the chooser}

\cli{monitor}, \cli{ps} and \cli{kill} with no argument all print the same thing --- the
scans currently alive on this machine, grouped from the \textsc{os} titles \Jtwo\ gives
its own processes, each with a short reference.\sidenote{On screen this arrives inside a
rounded, coloured panel whose width follows your terminal; the transcripts here show the
content with that frame stripped.}

\begin{terminal}
\begin{jcmd}
Jarvis ps
\end{jcmd}
\termsep
\begin{jout}
Running Jarvis scan tasks (2)

  REF   SCAN             PROCESSES   PIDS
  ----------------------------------------------------------------
  R1    EggBox_Dynesty           4   48120, 48127, 48131, 48140
  R2    NMSSM_02                 3   51002, 51009, 51015

Choose a task by its REF (for example R1):
  Jarvis monitor R1   view this task's live scan status
  Jarvis ps R1        list its control, workers, archiver, Redis
  Jarvis kill R1      terminate ALL processes belonging to this task
\end{jout}
\end{terminal}

\code{R1} and \code{R2} are assigned fresh on every listing, by scan name in
case-insensitive alphabetical order --- they are a convenience for the next command you
type, not identifiers to store in a script. Anywhere a reference is accepted the exact
scan name works too, which is what a script should use. With nothing running you get one
line, \code{No running Jarvis scan tasks.}, and exit code \code{0}: these commands report,
they do not judge.

\subsection{Step two: one scan}

Naming a reference turns the listing into a card for that scan alone --- its Redis
endpoint, its output directory, and one row per process with the \emph{role} that process
plays rather than its raw title:

\begin{terminal}
\begin{jcmd}
Jarvis ps R1
\end{jcmd}
\termsep
\begin{jout}
Runtime task - R1

  REF    R1
  SCAN   EggBox_Dynesty
  REDIS  127.0.0.1:6390   db 3
  OUTPUT /work/eggbox/outputs/EggBox_Dynesty

    PID   ROLE       DETAIL
  ----------------------------------------------------
  48120   Control    scan coordinator
  48127   Redis      127.0.0.1:6390
  48131   Archiver   writes DATABASE and exports
  48140   Worker 0   calculator / sampling worker
\end{jout}
\end{terminal}

The \code{REDIS} and \code{OUTPUT} rows are not read from your card --- they are read from
the runtime metadata the scan itself published into Redis, then checked against the
process you selected. If that cross-check fails, the command refuses rather than showing
you a card that might describe a different scan:

\begin{terminal}
\begin{jout}
Refusing unverified runtime: control PID does not match metadata.
\end{jout}
\end{terminal}

An unknown reference is a usage error (exit \code{2}) and reprints the chooser, so a typo
costs you nothing:

\begin{terminal}
\begin{jout}
Unknown running scan 'R9'; available: R1 (EggBox_Dynesty), R2 (NMSSM_02)
\end{jout}
\end{terminal}

\subsection{Ending one scan}

\begin{terminal}
\begin{jcmd}
Jarvis kill R1              # shows the card, then asks
Jarvis kill R1 --yes        # no prompt (batch scripts, cleanup hooks)
Jarvis kill R1 --no-force   # SIGTERM only; never escalate to SIGKILL
\end{jcmd}
\end{terminal}

\cli{kill} prints the same card, with the destructive consequence spelled out on it, then
asks; on confirmation it signals that scan's processes --- \textsc{sigterm}, a two-second
grace period, then \textsc{sigkill} for whatever is still standing --- and closes with a
tally and a verdict:

\begin{terminal}
\begin{jout}
DANGER: Jarvis kill R1 will terminate all 4 processes above.
Kill 4 Jarvis process(es)? [y/N] y
kill: SIGTERM=4 SIGKILL=0 gone=0 failed=0
All matched Jarvis processes terminated.
\end{jout}
\end{terminal}

Four details in that sequence are deliberate, and each one has bitten somebody:

\begin{enumerate}
  \item \textbf{Only the selected scan dies.} The other families on the machine are
        untouched, which is the whole reason the reference is mandatory rather than
        optional-with-a-default.
  \item \textbf{The order is dependents first}: Workers, then the file-operation
        helper, then the Archiver, then the control process, and Redis last. Killing the
        control process first would leave it a window in which to respawn the Workers
        you are trying to remove, and killing Redis first would strand every process
        that is mid-transaction against it.
  \item \textbf{Without a terminal, it refuses}. Piped or run from \cli{cron}, where
        nobody can answer the prompt, it declines rather than guessing --- the message
        names the flag it wants, \code{re-run with: Jarvis kill --yes}. Automation has
        to say so explicitly.
  \item \textbf{Anything left alive is reported, not glossed over}. Processes it could
        not signal are listed by \textsc{pid} and the exit code is \code{1}, so a
        cleanup script can tell ``all clear'' from ``still there, and it is not mine to
        kill''.
\end{enumerate}

\begin{jwarn}
\cli{kill} is the last resort, not the stop button. \textbf{Ctrl+C} on the running scan
shuts down cleanly and \emph{saves the checkpoint}; a signalled control process does
not, so the work since the last checkpoint barrier is gone
(Chapter~\ref{ch:checkpoint}). Reach for \cli{kill} when there is no longer a scan to
interrupt --- only its remains.
\end{jwarn}

\section{Discovery: the registries, without a scan}
\label{sec:cli-discovery}

\begin{terminal}
\begin{jcmd}
Jarvis portal man             # every registered file format
Jarvis portal man slha        # one format's manual
Jarvis portal io_task.yaml    # run a standalone Portal IO YAML
Jarvis operas list            # every registered operator
Jarvis operas info math.add   # one operator's signature and doc
\end{jcmd}
\end{terminal}

\cli{portal} and \cli{operas} are not curated subsets --- each forwards its whole argument
list to the neighbouring package's own \textsc{cli} (\cli{jportal} and \cli{jopera}
respectively) and returns its exit code. Whatever those tools can do, this spelling can do,
including their own \cli{-h}; nothing here has to be kept in sync as they grow. Missing
\code{Jarvis-Operas} is reported as such (exit \code{2}) rather than as an unknown
command. \cli{Jarvis portal man} prints what \Portal\ can currently read and write:

\begin{terminal}
\begin{jcmd}
Jarvis portal man
\end{jcmd}
\termsep
\begin{jout}
Jarvis-Portal manual

Supported formats:
  input: CSV, DAT, JSON, SLHA, TSV, Wolfram
  output: CSV, DAT, JSON, SLHA, TSV, Wolfram, xSLHA
\end{jout}
\end{terminal}

The reason to prefer these over the documentation --- including this book --- is that
they interrogate the registries in \emph{your} environment. \cli{operas list} names the
operators actually importable here, so a missing extra or a stale plugin version shows
up as an absent name rather than as a puzzling \yamlkey{expression} failure halfway
through a scan. When a card names something these commands do not list, the card is
wrong about this machine.

\section{Legacy spellings}
\label{sec:cli-legacy}

\Jarvis\ V1 was driven by a bare card path plus a mode flag, and those spellings still
work. The rule is mechanical: when the first token is neither a known subcommand nor a
flag, \Jtwo\ treats it as a path and lets a trailing legacy flag pick the intent.

\begin{itemize}
  \item \cli{Jarvis <task>.yaml} routes to \cli{run}.
  \item \cli{Jarvis <task>.yaml --resume} routes to \cli{run --resume}.
  \item \cli{Jarvis <task>.yaml --check-modules} routes to \cli{check}.
  \item \cli{Jarvis <task>.yaml --validate} routes to \cli{validate}.
  \item \cli{Jarvis <task>.yaml --convert} routes to \cli{convert}.
  \item \cli{Jarvis --monitor} routes to \cli{monitor} (\S\ref{sec:cli-ps-kill}).
  \item \cli{Jarvis <scene>.yaml --plot} routes to \cli{plot}, with a deprecation
        warning.
\end{itemize}

That list is complete, and \cli{Jarvis -h} prints exactly these in their own
\emph{Legacy options} panel, kept apart from the real options so the two are never
confused for each other.

Prefer the subcommand forms in anything you save. They survive shell quoting, they read
correctly in a batch script somebody else inherits, and they cannot collide: naming a
subcommand \emph{and} a conflicting legacy flag in one line is rejected as a usage error
rather than resolved by precedence.

\section{Exit codes}
\label{sec:exit-codes}

\clih{Jarvis} reports the outcome of every command through its exit status, listed in
Table~\ref{tab:exit-codes}.

\begin{table}[ht]
  \footnotesize
  \begingroup
  \setlength{\tabcolsep}{0pt}
  \begin{tabular}{@{}p{0.34\linewidth}p{0.66\linewidth}@{}}
    \toprule
    \textbf{Code} & \textbf{Meaning} \\
    \midrule
    0   & every sample succeeded \\
    1   & any sample failed (including partial failure), or a runtime/IO error \\
    2   & usage or configuration error (bad flags, missing/invalid card) \\
    130 & interrupted (Ctrl+C) \\
    \bottomrule
  \end{tabular}
  \endgroup
  \caption{\clih{Jarvis} exit codes --- contract-stable for scripting.}
  \label{tab:exit-codes}
\end{table}

These four are contract-stable, and the run codes are computed from the \code{status}
field of the closing \code{RunOutcome} line: \code{success} gives \code{0};
\code{partial\_failure}, \code{failed}, and \code{error} all give \code{1};
\code{interrupted} gives \code{130}. Reporting one number for three different kinds of
trouble is the deliberate part.

So is treating \emph{partial} failure as failure. A scan where nineteen of twenty points
completed exits \code{1}, because the alternative --- calling it success --- is how a
pipeline quietly ships a dataset with a hole in it. Failed samples are archived as failed
records, so a workflow that legitimately tolerates ``some points crashed'' should read
the counters in \file{run\_summary.json} rather than widening its notion of a good exit
code.

\section{Environment variables}
\label{sec:cli-env}

Table~\ref{tab:cli-env} lists the environment variables the \textsc{cli} reads.

\begin{table}[ht]
  \footnotesize
  \begingroup
  \setlength{\tabcolsep}{0pt}
  \begin{tabular}{@{}p{0.34\linewidth}p{0.66\linewidth}@{}}
    \toprule
    \textbf{Variable} & \textbf{Effect} \\
    \midrule
    \code{JARVIS\_HEP\_TASK\_ROOT} / \code{JHEP\_TASK\_ROOT} & force the project root
      (Section~\ref{sec:projects}) \\
    \code{JARVIS\_PROJECT\_FETCH\_KEY} & decryption key for restricted fetches
      (Ch.~\ref{ch:project-cli}) \\
    \code{JARVIS\_OFFICIAL\_LIBRARY\_INDEX\_URL} & override the catalog index URL
      (\code{https://} or \code{file://}) \\
    \code{JARVIS\_OFFICIAL\_LIBRARY\_TIMEOUT\_SEC} & catalog fetch timeout \\
    \bottomrule
  \end{tabular}
  \endgroup
  \caption{Environment variables read by the CLI.}
  \label{tab:cli-env}
\end{table}

\begin{jtip}
\textbf{The habit, in three lines.} Gate the card in \textsc{ci} with
\cli{Jarvis validate --strict}; after every edit that touches a path, a program, or an
expression, spend a minute on \cli{Jarvis check}; then \cli{Jarvis run --strict} once,
and watch it from a second terminal with \cli{Jarvis monitor} followed by the reference
it hands you. Every rung you skip is one you will climb later, at the price of the rung
above it.
\end{jtip}
